\documentclass[10pt,a4paper]{amsart}
\usepackage[margin=19mm]{geometry}
\usepackage{amsmath,amssymb,mathtools}
\usepackage[scr=rsfs,cal=euler]{mathalfa}
\usepackage{xfrac}
\usepackage[unicode,hidelinks,bookmarks=false]{hyperref}
\newcommand{\ie}{i.e.\ }
\newcommand{\cf}{cf.\ }
\newcommand{\resp}{respectively\ }
\newcommand{\Subsection}{\subsection}
\providecommand{\nobreakdash}{\nobreak\hbox{-}}
\setlength{\emergencystretch}{2em}
\multlinegap0pt
\usepackage{amssymb}
\usepackage{enumerate}
\usepackage{xcolor}
\usepackage{tikz}
\newtheorem{proposition}{Proposition}
\title{Bellis strong stable sets on infinite hyperbolic surfaces}
\author{Sergi Burniol Clotet, Fran\c{c}oise Dal'Bo, Sergio Herrero Vila}
\date{Source: arXiv:2604.02649v1 (CC BY 4.0)}
\begin{document}
\maketitle

\noindent\emph{Source and licence.} Bellis strong stable sets on infinite hyperbolic surfaces. Version arXiv:2604.02649v1 (CC BY 4.0), distributed by arXiv under the Creative Commons Attribution 4.0 International licence, http://creativecommons.org/licenses/by/4.0/, as recorded in the arXiv licence metadata for this exact version and shown on the abstract page https://arxiv.org/abs/2604.02649v1. Full licence notice: \texttt{licenses/arxiv-2604.02649-CC-BY-4.0.txt}.\par\smallskip

\noindent\textit{Editorial scope.} Counted unit: the article's proposition prop4 (source0.tex lines 1785--1799). The article's complete original proof of it is delivered with it (source0.tex lines 1801--1934, one \texttt{proof} environment of 134 lines). No mathematics is written, repaired or re-derived: the statement and the proof are the article's own lines, in the article's own notation, and no retained line is substituted or re-flowed.  No further source-local context is needed: every cross-reference of every retained line resolves inside the two delivered spans.  Nothing was omitted from any retained span, so the package carries no omission record.  No retained line contains a citation, so the wrapper carries no bibliography and no bibliography entry.  No retained line contains a figure environment, an image-inclusion command or drawing code, so no image is delivered and the package declares no asset dependency.  Editorial formatting only, in addition: an AMS \texttt{amsart} preamble that restates the article's own declarations and macros used by the retained lines, namely the environments \texttt{proposition} on separate counters, together with the standard packages those lines need.  The retained environments print in this document as Proposition 1 in order of appearance, whereas the article prints the same items with the numbering of its own full document.  The complete native source of the article as distributed in the arXiv e-print for this exact version is bundled unchanged as \texttt{sources/197753/source0.tex}.
\begin{proposition}\label{prop4}
	Let $\tilde u_{0}\in T^{1}\mathbb H^{2}$ such that
	\[
	\tilde u_{0}(0)=i
	\qquad\text{and}\qquad
	\tilde u_{0}(+\infty)=\infty.
	\]
	Let $(\tilde v_{t})_{t\geq 0}$ be a family of vectors in $T^{1}\mathbb H^{2}$.
	Then
	\[
	\lim\limits_{t\to+\infty} d_{1}(\tilde u_{0},\tilde v_{t})=0
	\qquad\Longleftrightarrow\qquad
	\lim\limits_{t\to+\infty} d_{2}(\tilde u_{0},\tilde v_{t})=0.
	\]
\end{proposition}

\begin{proof}

For $t\geq 0$, set
\begin{itemize}
	
	\item$ \alpha_{t}=\alpha_{\tilde u_{0},\tilde v_{t}},
	\qquad
	\beta_{t}=\beta_{\tilde u_{0},\tilde v_{t}},
	$
	\item $\xi_{t}=C_{\tilde u_{0},\tilde v_{t}}(+\infty),
	\qquad
	\xi_{t}'=\tilde v_{t}(+\infty).
	$
\end{itemize}
\medskip

\noindent
$\left( \Longrightarrow \right)$ Suppose $
\lim\limits_{t\to+\infty} d_{1}(\tilde u_{0},\tilde v_{t})=0.
$

Since $
\lim\limits_{t\to+\infty}\tilde v_{t}(0)=i
\text{ and }
\lim\limits_{t\to+\infty}\tilde v_{t}(1)=e\,i,
$
we have $
\lim\limits_{t\to+\infty}\xi_{t}'=\infty.
$

Suppose by contradiction that there exist $\varepsilon>0$ and a sequence $(t_{n})_{n\geq 0}\subset\mathbb R_{+}$
converging to $+\infty$ such that
\begin{equation}\label{eq1}
	\forall n\geq 0,\qquad |\alpha_{t_{n}}-\beta_{t_{n}}|\geq \varepsilon.
\end{equation}

Up to a subsequence, we can suppose $
\lim\limits_{n\to+\infty}\xi_{t_{n}}=\xi.
$

Let $\tilde w\in T^{1}\mathbb H^{2}$ such that
$
\tilde w(0)=i
\text{ and }
\tilde w(+\infty)=\xi,
$

we have
$
\lim\limits_{n\to+\infty} C_{\tilde u_{0},\tilde v_{t_{n}}}=\tilde w,
\lim\limits_{n\to+\infty}g_{s_n}C_{\tilde u_0,\tilde v_{t_n}}=\tilde w, \text{ with } s_n=d(i,\tilde v_{t_n}(0))
\text{ and }
\lim\limits_{n\to+\infty}\tilde v_{t_{n}}=\tilde u_{0}.
$

It follows that
\[
\lim\limits_{n\to+\infty}\alpha_{t_{n}}=\widehat{(\tilde w,\tilde u_{0})}
\qquad\text{and}\qquad
\lim\limits_{n\to+\infty}\beta_{t_{n}}=\widehat{(\tilde w,\tilde u_{0})},
\]
hence
$
\lim\limits_{n\to+\infty}|\alpha_{t_{n}}-\beta_{t_{n}}|=0,
$
which contradicts (\ref{eq1}).

\medskip

\noindent
$\left( \Longleftarrow \right)$ Suppose
$
\lim\limits_{t\to+\infty} d_{2}(\tilde u_{0},\tilde v_{t})=0.
$

Suppose by contradiction that there exist $\varepsilon>0$ and $(t_{n})_{n\geq 0}\subset\mathbb R_{+}$
converging to $+\infty$ such that
\begin{equation}\label{eq2}
	\forall t_{n},\qquad d\bigl(\tilde u_{0}(1),\tilde v_{t_{n}}(1)\bigr)\geq \varepsilon.
\end{equation}

We can suppose
$
\lim\limits_{n\to+\infty}\xi_{t_{n}}=\xi
\text{ and }
\lim\limits_{n\to+\infty}\xi_{t_{n}}'=\xi'.
$

Let $\tilde w\in T^{1}\mathbb H^{2}$ with
$
\tilde w(0)=i
\text{ and }
\tilde w(+\infty)=\xi,
$
and $\tilde w'\in T^{1}\mathbb H^{2}$ with
$
\tilde w'(0)=i
\text{ and } 
\tilde w'(+\infty)=\xi'.
$

Clearly,
$
\lim\limits_{n\to+\infty}\tilde v_{t_{n}}=\tilde w'
\text{ and } 
\lim\limits_{n\to+\infty} C_{\tilde u_{0},\tilde v_{t_{n}}}=\tilde w.
$

It follows that
$
\lim\limits_{n\to+\infty}\alpha_{t_{n}}=\widehat{(\tilde w,\tilde u_{0})}
\text{ and } 
\lim\limits_{n\to+\infty}\beta_{t_{n}}=\widehat{(\tilde w,\tilde w')}.
$ 

Since
$
\lim\limits_{n\to+\infty}|\alpha_{t_{n}}-\beta_{t_{n}}|=0,
$
we obtain
$
\tilde w'=\tilde u_{0}.
$

We deduce that the geodesic segment
$
[\tilde v_{t_{n}}(0),\tilde v_{t_{n}}(1)]
$
converges to the vertical arc
$
[i,e\,i],
$
which contradicts (\ref{eq2}).
\end{proof}

\end{document}
